\chapter{Casos de uso - Segunda iteración}
\label{chap:anexo_casos_uso_it2}

Este anexo presenta la especificación detallada de los casos de uso desarrollados durante la segunda iteración del sistema, enfocados en la refactorización arquitectónica de los módulos independientes para facilitar futuras integraciones.

\section{CU-01 - Generación de datasets para entrenamiento ML}
\label{sec:cu01_it2_datasets_ml}

\begin{table}[H]
\centering
\resizebox{\textwidth}{!}{
\rowcolors{2}{white}{udcgray!25}
\begin{tabular}{|p{3cm}|p{12cm}|}
\hline
\rowcolor{udcpink!25}
\textbf{Campo} & \textbf{Descripción} \\
\hline
\textbf{Identificador} & CU-01 \\
\hline
\textbf{Nombre} & Generación de datasets mediante módulos independientes \\
\hline
\textbf{Actores} & Investigador BCI (principal), Analista de datos (secundario) \\
\hline
\textbf{Descripción} & El sistema facilita la generación de datasets mediante dos módulos independientes que pueden ejecutarse simultáneamente para capturar datos sincronizables offline \\
\hline
\textbf{Precondiciones} & 
\begin{minipage}[t]{11cm}
\begin{itemize}
    \item Ambos módulos están funcionando correctamente
    \item Los dispositivos de entrada y EEG están configurados
    \item Existe al menos un formato de exportación implementado
    \item Los módulos permiten operación independiente
\end{itemize}
\end{minipage} \\
\hline
\textbf{Postcondiciones} & 
\begin{minipage}[t]{11cm}
\begin{itemize}
    \item Cada módulo genera archivos independientes con timestamps precisos
    \item Los datos están estructurados para sincronización offline posterior
    \item Los archivos facilitan la correlación temporal entre modalidades
    \item La arquitectura modular facilita el mantenimiento del sistema
\end{itemize}
\end{minipage} \\
\hline
\end{tabular}
}
\caption{Especificación del caso de uso CU-01 - Segunda iteración}
\label{tab:cu01_it2_especificacion}
\end{table}

\subsection{Flujo principal}
\label{subsec:cu01_it2_flujo_principal}

\begin{table}[H]
\centering
\resizebox{\textwidth}{!}{
\rowcolors{2}{white}{udcgray!25}
\begin{tabular}{|c|p{13cm}|}
\hline
\rowcolor{udcpink!25}
\textbf{Paso} & \textbf{Acción} \\
\hline
1 & El investigador BCI inicia ambos módulos de forma independiente \\
\hline
2 & El investigador configura parámetros específicos en cada módulo por separado \\
\hline
3 & Se coordinan manualmente los tiempos de inicio de captura en ambos módulos \\
\hline
4 & Cada módulo opera independientemente registrando datos con timestamps \\
\hline
5 & El usuario del sistema realiza acciones experimentales \\
\hline
6 & Los módulos registran eventos con precisión temporal en archivos separados \\
\hline
7 & Se coordinan manualmente los tiempos de finalización de captura \\
\hline
8 & Cada módulo genera archivos independientes listos para análisis offline \\
\hline
9 & Los timestamps permiten sincronización posterior mediante herramientas externas \\
\hline
\end{tabular}
}
\caption{Flujo principal del caso de uso CU-01 - Segunda iteración}
\label{tab:cu01_it2_flujo_principal}
\end{table}

\subsection{Flujos alternativos}
\label{subsec:cu01_it2_flujos_alternativos}

\begin{table}[H]
\centering
\resizebox{\textwidth}{!}{
\rowcolors{2}{white}{udcgray!25}
\begin{tabular}{|p{3cm}|p{11cm}|}
\hline
\rowcolor{udcpink!25}
\textbf{Flujo} & \textbf{Descripción} \\
\hline
\textbf{FA-01: Inicio escalonado} & 
\begin{minipage}[t]{10.5cm}
\begin{itemize}
    \item Los módulos se inician en momentos diferentes según protocolo
    \item Los timestamps independientes mantienen trazabilidad temporal
    \item La sincronización offline compensa diferencias de inicio
    \item Permite protocolos experimentales que requieren preparación secuencial
\end{itemize}
\end{minipage} \\
\hline
\textbf{FA-02: Captura independiente extendida} & 
\begin{minipage}[t]{10.5cm}
\begin{itemize}
    \item Un módulo puede continuar capturando después del otro
    \item Los archivos reflejan la duración real de cada captura
    \item Útil para análisis de períodos pre/post experimentales
    \item Mantiene integridad de datos en cada modalidad
\end{itemize}
\end{minipage} \\
\hline
\end{tabular}
}
\caption{Flujos alternativos del caso de uso CU-01 - Segunda iteración}
\label{tab:cu01_it2_flujos_alternativos}
\end{table}

\subsection{Flujos de excepción}
\label{subsec:cu01_it2_flujos_excepcion}

\begin{table}[H]
\centering
\resizebox{\textwidth}{!}{
\rowcolors{2}{white}{udcgray!25}
\begin{tabular}{|p{3cm}|p{11cm}|}
\hline
\rowcolor{udcpink!25}
\textbf{Excepción} & \textbf{Manejo} \\
\hline
\textbf{FE-01: Fallo de un módulo independiente} & 
\begin{minipage}[t]{10.5cm}
\begin{itemize}
    \item El otro módulo continúa funcionando sin interrupción
    \item Los datos del módulo funcional se preservan completamente
    \item El fallo no afecta la integridad del módulo operativo
    \item Los archivos parciales mantienen validez para análisis
\end{itemize}
\end{minipage} \\
\hline
\textbf{FE-02: Desincronización temporal} & 
\begin{minipage}[t]{10.5cm}
\begin{itemize}
    \item Los módulos mantienen sus timestamps internos independientes
    \item La desincronización se detecta durante análisis offline
    \item Herramientas externas pueden compensar diferencias temporales
    \item La arquitectura modular facilita diagnóstico de problemas
\end{itemize}
\end{minipage} \\
\hline
\end{tabular}
}
\caption{Flujos de excepción del caso de uso CU-01 - Segunda iteración}
\label{tab:cu01_it2_flujos_excepcion}
\end{table}

\section{CU-03 - Desarrollo de aplicaciones BCI personalizadas (fundamentos)}
\label{sec:cu03_it2_desarrollo_bci}

\begin{table}[H]
\centering
\resizebox{\textwidth}{!}{
\rowcolors{2}{white}{udcgray!25}
\begin{tabular}{|p{3cm}|p{12cm}|}
\hline
\rowcolor{udcpink!25}
\textbf{Campo} & \textbf{Descripción} \\
\hline
\textbf{Identificador} & CU-03 \\
\hline
\textbf{Nombre} & Fundamentos para desarrollo de extensiones modulares \\
\hline
\textbf{Actores} & Desarrollador de aplicaciones BCI (principal), Integrador de sistemas (secundario) \\
\hline
\textbf{Descripción} & La refactorización arquitectónica proporciona una base modular que facilita el desarrollo futuro de extensiones y aplicaciones personalizadas \\
\hline
\textbf{Precondiciones} & 
\begin{minipage}[t]{11cm}
\begin{itemize}
    \item Los módulos están completamente refactorizados según principios SOLID
    \item Los componentes están separados y bien definidos
    \item Las interfaces proporcionan acceso controlado a funcionalidades
    \item La documentación técnica está actualizada
\end{itemize}
\end{minipage} \\
\hline
\textbf{Postcondiciones} & 
\begin{minipage}[t]{11cm}
\begin{itemize}
    \item Las extensiones pueden integrarse con mínimo impacto en código base
    \item Las interfaces modulares facilitan desarrollo de terceros
    \item La arquitectura soporta evolución independiente de componentes
    \item Se facilita el mantenimiento y extensión del sistema
\end{itemize}
\end{minipage} \\
\hline
\end{tabular}
}
\caption{Especificación del caso de uso CU-03 - Segunda iteración}
\label{tab:cu03_it2_especificacion}
\end{table}

\subsection{Flujo principal}
\label{subsec:cu03_it2_flujo_principal}

\begin{table}[H]
\centering
\resizebox{\textwidth}{!}{
\rowcolors{2}{white}{udcgray!25}
\begin{tabular}{|c|p{13cm}|}
\hline
\rowcolor{udcpink!25}
\textbf{Paso} & \textbf{Acción} \\
\hline
1 & El desarrollador analiza la arquitectura refactorizada y componentes disponibles \\
\hline
2 & Identifica puntos de extensión en la arquitectura modular \\
\hline
3 & Diseña extensión siguiendo los principios arquitectónicos establecidos \\
\hline
4 & Implementa interfaces requeridas según la funcionalidad deseada \\
\hline
5 & Desarrolla lógica personalizada aprovechando componentes existentes \\
\hline
6 & Integra extensión mediante modificaciones mínimas del código base \\
\hline
7 & Configura comunicación con componentes apropriados \\
\hline
8 & Valida que la extensión respeta la separación de responsabilidades \\
\hline
9 & Prueba integración verificando que no afecta funcionalidad base \\
\hline
10 & Documenta extensión según estándares del proyecto refactorizado \\
\hline
\end{tabular}
}
\caption{Flujo principal del caso de uso CU-03 - Segunda iteración}
\label{tab:cu03_it2_flujo_principal}
\end{table}

\subsection{Flujos alternativos}
\label{subsec:cu03_it2_flujos_alternativos}

\begin{table}[H]
\centering
\resizebox{\textwidth}{!}{
\rowcolors{2}{white}{udcgray!25}
\begin{tabular}{|p{3cm}|p{11cm}|}
\hline
\rowcolor{udcpink!25}
\textbf{Flujo} & \textbf{Descripción} \\
\hline
\textbf{FA-01: Extensión de colector personalizado} & 
\begin{minipage}[t]{10.5cm}
\begin{itemize}
    \item Desarrollador implementa nueva interfaz de colector
    \item Utiliza la arquitectura de gestores para integración
    \item El nuevo colector se registra en el sistema modular
    \item Sistema incluye funcionalidad sin modificar código base
\end{itemize}
\end{minipage} \\
\hline
\textbf{FA-02: Visualizador especializado} & 
\begin{minipage}[t]{10.5cm}
\begin{itemize}
    \item Desarrollador extiende arquitectura de visualización
    \item Implementa comunicación con componentes de datos
    \item Registra visualizador en sistema de gestión correspondiente
    \item Integración aprovecha separación de responsabilidades existente
\end{itemize}
\end{minipage} \\
\hline
\end{tabular}
}
\caption{Flujos alternativos del caso de uso CU-03 - Segunda iteración}
\label{tab:cu03_it2_flujos_alternativos}
\end{table}

\subsection{Flujos de excepción}
\label{subsec:cu03_it2_flujos_excepcion}

\begin{table}[H]
\centering
\resizebox{\textwidth}{!}{
\rowcolors{2}{white}{udcgray!25}
\begin{tabular}{|p{3cm}|p{11cm}|}
\hline
\rowcolor{udcpink!25}
\textbf{Excepción} & \textbf{Manejo} \\
\hline
\textbf{FE-01: Componentes no disponibles} & 
\begin{minipage}[t]{10.5cm}
\begin{itemize}
    \item Extensión requiere componentes no inicializados
    \item Sistema muestra error específico en terminal
    \item La arquitectura modular aísla el fallo
    \item Resto del sistema continúa funcionando normalmente
\end{itemize}
\end{minipage} \\
\hline
\textbf{FE-02: Error durante ejecución de extensión} & 
\begin{minipage}[t]{10.5cm}
\begin{itemize}
    \item Excepción no controlada en código de extensión
    \item La separación arquitectónica contiene el error
    \item Sistema continúa con degradación elegante
    \item Se pueden mostrar logs de depuración específicos
\end{itemize}
\end{minipage} \\
\hline
\end{tabular}
}
\caption{Flujos de excepción del caso de uso CU-03 - Segunda iteración}
\label{tab:cu03_it2_flujos_excepcion}
\end{table}

\section{CU-06 - Preparación para exportación de datos}
\label{sec:cu06_it2_exportacion_datos}

\begin{table}[H]
\centering
\resizebox{\textwidth}{!}{
\rowcolors{2}{white}{udcgray!25}
\begin{tabular}{|p{3cm}|p{12cm}|}
\hline
\rowcolor{udcpink!25}
\textbf{Campo} & \textbf{Descripción} \\
\hline
\textbf{Identificador} & CU-06 \\
\hline
\textbf{Nombre} & Preparación de datos para análisis posterior con herramientas externas \\
\hline
\textbf{Actores} & Analista de datos (principal), Investigador BCI (secundario) \\
\hline
\textbf{Descripción} & Los módulos refactorizados facilitan la generación de archivos estructurados que pueden ser utilizados posteriormente con herramientas de análisis externas \\
\hline
\textbf{Precondiciones} & 
\begin{minipage}[t]{11cm}
\begin{itemize}
    \item Los módulos están refactorizados y funcionando correctamente
    \item Los sistemas de exportación soportan formato CSV
    \item Los datos se escriben sobre la marcha durante la captura
    \item La arquitectura modular facilita futuras mejoras
\end{itemize}
\end{minipage} \\
\hline
\textbf{Postcondiciones} & 
\begin{minipage}[t]{11cm}
\begin{itemize}
    \item Los datos están disponibles en formato CSV estructurado
    \item Los archivos conservan timestamps precisos para análisis temporal
    \item Los metadatos facilitan la interpretación de los datos
    \item Los archivos están listos para importación en herramientas externas
\end{itemize}
\end{minipage} \\
\hline
\end{tabular}
}
\caption{Especificación del caso de uso CU-06 - Segunda iteración}
\label{tab:cu06_it2_especificacion}
\end{table}

\subsection{Flujo principal}
\label{subsec:cu06_it2_flujo_principal}

\begin{table}[H]
\centering
\resizebox{\textwidth}{!}{
\rowcolors{2}{white}{udcgray!25}
\begin{tabular}{|c|p{13cm}|}
\hline
\rowcolor{udcpink!25}
\textbf{Paso} & \textbf{Acción} \\
\hline
1 & El analista identifica necesidad de datos para análisis con herramientas externas \\
\hline
2 & Configura captura en módulos independientes para generar archivos CSV \\
\hline
3 & Inicia captura con exportación activa en ambos módulos \\
\hline
4 & Los módulos escriben datos estructurados sobre la marcha \\
\hline
5 & Se incluyen metadatos relevantes para interpretación posterior \\
\hline
6 & Cada módulo genera archivos CSV con timestamps precisos \\
\hline
7 & Completa captura con archivos CSV listos para uso \\
\hline
8 & Verifica estructura y contenido de archivos generados \\
\hline
9 & Los archivos están disponibles para herramientas de análisis externas \\
\hline
10 & El analista puede proceder con análisis offline utilizando los datos \\
\hline
\end{tabular}
}
\caption{Flujo principal del caso de uso CU-06 - Segunda iteración}
\label{tab:cu06_it2_flujo_principal}
\end{table}

\subsection{Flujos de excepción}
\label{subsec:cu06_it2_flujos_excepcion}

\begin{table}[H]
\centering
\resizebox{\textwidth}{!}{
\rowcolors{2}{white}{udcgray!25}
\begin{tabular}{|p{3cm}|p{11cm}|}
\hline
\rowcolor{udcpink!25}
\textbf{Excepción} & \textbf{Manejo} \\
\hline
\textbf{FE-01: Error durante escritura} & 
\begin{minipage}[t]{10.5cm}
\begin{itemize}
    \item Error no controlado durante proceso de escritura de archivos
    \item El módulo afectado se detiene y puede mostrar mensaje de debug
    \item Los datos escritos hasta el momento se preservan
    \item Los metadatos de finalización no se generan para ese módulo
\end{itemize}
\end{minipage} \\
\hline
\end{tabular}
}
\caption{Flujos de excepción del caso de uso CU-06 - Segunda iteración}
\label{tab:cu06_it2_flujos_excepcion}
\end{table}